% =====================================================================
% Algoritmos de Grafo e Ancoragem das Respostas
% Responsável: Pessoa 5
%
% O que cobrir:
% - Algoritmos do Neo4j GDS usados (PageRank, Louvain/Leiden, Shortest Path) e para que cada um serve no contexto passado ao LLM.
% - Exemplo concreto: entidades mais centrais do grafo pelo PageRank, comunidades encontradas.
% - Grounding: como o LLM é obrigado a citar os caminhos do subgrafo e os trechos (evidências).
% - Como o nível de confiança da resposta é calculado.
%
% Dicas: use \rot{Pessoa} para rótulos e \rel{CRIOU} para relações.
% Cite com \cite{chave} (chaves em referencias.bib). Ao terminar,
% apague o \pendente abaixo.
% =====================================================================

\section{Algoritmos de Grafo e Ancoragem das Respostas}
\label{sec:grounding}

\pendente{algoritmos de grafo e ancoragem das respostas}
